% Part: first-order-logic
% Chapter: tableaux
% Section: provability-consistency

\documentclass[../../../include/open-logic-section]{subfiles}

\begin{document}

\iftag{FOL}
      {\olfileid{fol}{tab}{prv}}
      {\olfileid{pl}{tab}{prv}}

\olsection{\usetoken{S}{derivability} او سازګاري}

اوس د~!!{derivability} اړيکې يو شمېر خاصيتونه ثابتوو۔ دوئ په
خپلواک ډول هم په زړه پورې دي، خو هر يو به د بشپړتيا قضیې په ثبوت
کښې ونډه ولري۔

\begin{prop}\ollabel{prop:provability-contr}
  که $\Gamma \Proves !A$ او $\Gamma \cup \{!A\}$ ناسازګار وي،
  نو $\Gamma$ ناسازګار دے۔
\end{prop}

\begin{proof}
داسې متناهي
$\Gamma_0 = \{!B_1, \dots, !B_n\}$ او
$\Gamma_1 =\{!C_1, \dots, !C_m\} \subseteq \Gamma$ شته چې
  \begin{align*}
    \{\sFmla{\False}{!A}, &
    \sFmla{\True}{!B_1}, \dots, \sFmla{\True}{!B_n}\} \\
    \{\sFmla{\True}{!A}, &
    \sFmla{\True}{!C_1}, \dots, \sFmla{\True}{!C_m}\}
  \end{align*}
  تړلي !!{tableau}s لري۔ پر~$!A$ د~\Cut{} قاعدې په کارولو دا په
  يوه تړلي !!{tableau} کښې يو ځاے کولے شو چې د
  $\Gamma_0 \cup \Gamma_1$ ناسازګاري ښيي۔ ځکه
  $\Gamma_0 \subseteq \Gamma$ او~$\Gamma_1 \subseteq \Gamma$،
  نو $\Gamma_0 \cup \Gamma_1 \subseteq \Gamma$؛ له دې امله
  $\Gamma$ ناسازګار دے۔
\textbf{د ژباړې د نسخې سمون (OLTAB-005):} سرچينه د دويم متناهي
سټ په تعريف کښې لړ په اېن شاخص پاے ته رسوي، خو په تابلو کښې ئې
وروستے غړے د اېم شاخص لري؛ دلته د لړ وروستے شاخص يو شان شوے دے۔
\end{proof}

\begin{prop}
\ollabel{prop:prov-incons}
$\Gamma \Proves !A$ هله او يوازې هله چې
$\Gamma \cup \{\lnot !A\}$ ناسازګار وي۔
\end{prop}

\begin{proof}
لومړے فرض کړئ $\Gamma \Proves !A$، يعنې د لاندې دپاره تړلے
!!{tableau} شته:
\[
\{\sFmla{\False}{!A},
\sFmla{\True}{!B_1}, \dots, \sFmla{\True}{!B_n}\}
\]
د~$\TRule{\True}{\lnot}$ قاعدې په کارولو دا د لاندې دپاره په تړلي
!!{tableau} اړولے شو:
\[
\{\sFmla{\True}{\lnot !A},
\sFmla{\True}{!B_1}, \dots, \sFmla{\True}{!B_n}\}.
\]

بل لوري ته، که د وروستي سټ دپاره تړلے !!{tableau} وي، د لومړي سټ
په تړلي !!{tableau} ئې اړولے شو: هغه هر فارمول ړنګوو چې پر لومړي
فرض~$\sFmla{\True}{\lnot !A}$ د~\TRule{\True}{\lnot} له کارونې
راوتلے وي، خپله دغه فرض هم ړنګوو، او فرض
$\sFmla{\False}{!A}$ ورزياتوو۔ ځکه که کومه څانګه مخکښې د
\TRule{\True}{\lnot} د هغې پايلې له امله تړلې وه چې پر
$\sFmla{\True}{\lnot !A}$ پلې شوې، يعنې
$\sFmla{\False}{!A}$، نو په نوي !!{tableau} کښې اړونده څانګه هم
تړلې ده۔ که په زوړ تابلو کښې کومه څانګه د فرض
$\sFmla{\True}{\lnot !A}$ او~$\sFmla{\False}{\lnot !A}$ دواړو له
امله تړلې وه، پر~$\sFmla{\False}{\lnot !A}$ د
$\TRule{\False}{\lnot}$ په پلي کولو او د
$\sFmla{\True}{!A}$ په اخستلو ئې په تړلې څانګه اړولے شو۔ دا څانګه
ځکه تړل کېږي چې $\sFmla{\False}{!A}$ مو د فرض په توګه ورزيات کړے
دے۔
\end{proof}

\begin{prob}
ثابت کړئ چې $\Gamma \Proves \lnot !A$ هله او يوازې هله چې
$\Gamma \cup \{!A\}$ ناسازګار وي۔
\end{prob}

\begin{prop}\ollabel{prop:explicit-inc}
  که $\Gamma \Proves !A$ او~$\lnot !A \in \Gamma$ وي، نو
  $\Gamma$ ناسازګار دے۔
\end{prop}

\begin{proof}
  فرض کړئ $\Gamma \Proves !A$ او~$\lnot !A \in \Gamma$۔ بيا داسې
  $!B_1$، \dots، $!B_n \in \Gamma$ شته چې
  \[
  \{\sFmla{\False}{!A}, \sFmla{\True}{!B_1}, \dots, \sFmla{\True}{!B_n}\}
  \]
  تړلے تابلو لري۔ فرض~\sFmla{\False}{!A} په
  \sFmla{\True}{\lnot !A} بدل کړئ، او د فرضونو وروسته هغه
  \sFmla{\False}{!A} پايله ورننباسئ چې پر
  \sFmla{\True}{\lnot !A} د~\TRule{\True}{\lnot} له کارونې
  راځي۔ په زوړ !!{tableau} کښې هره !!{sentence} چې د~$1$ کرښې په
  حواله توجيه شوې وه، اوس د~$n+1$ کرښې په حواله توجيه کېږي۔ نو که
  زوړ !!{tableau} تړلے و، نوے !!{tableau} هم تړلے دے۔ دا د~$\Gamma$ ناسازګاري
  ښيي؛ ځکه ټول فرضونه په~$\Gamma$ کښې دي۔
  \textbf{د ژباړې د نسخې سمون (OLTAB-006):} سرچينه د «داسې چې»
  وروسته يو بې دليله تېک شاتګی ږدي؛ دلته لرې شوے دے۔
  \textbf{د ژباړې د نسخې سمون (OLTAB-007):} سرچينه وايي د رښتونې
  نفي قاعده پر دروغې جملې پلې کېږي؛ قاعده په حقيقت کښې پر رښتونې
  نفي شوې جملې پلې کېږي او دروغ مثبت فارمول راوباسي، لکه د همدې
  ثبوت په راتلونکې جمله کښې۔
\end{proof}

\begin{prop}\ollabel{prop:provability-exhaustive}
  که $\Gamma \cup \{!A\}$ او~$\Gamma \cup \{\lnot !A\}$ دواړه
  ناسازګار وي، نو $\Gamma$ ناسازګار دے۔
\end{prop}

\begin{proof}
  که داسې $!B_1$، \dots، $!B_n \in \Gamma$ او
  $!C_1$، \dots، $!C_m \in \Gamma$ وي چې
  \begin{align*}
    \{\sFmla{\True}{!A}, &
    \sFmla{\True}{!B_1}, \dots, \sFmla{\True}{!B_n}\} \text{ او}\\
    \{\sFmla{\True}{\lnot !A}, &
    \sFmla{\True}{!C_1}, \dots, \sFmla{\True}{!C_m}\}
  \end{align*}
  دواړه تړلي !!{tableau}s ولري، يو داسې ګډ !!{tableau} جوړولے شو
  چې د~$\Gamma$ ناسازګاري ښيي۔ د فرضونو په توګه
  $\sFmla{\True}{!B_1}$، \dots، $\sFmla{\True}{!B_n}$ له
  $\sFmla{\True}{!C_1}$، \dots، $\sFmla{\True}{!C_m}$ سره يو
  ځاے کاروو، او بيا د~\Cut{} قاعده پلې کوو۔ دا دوه څانګې ورکوي:
  يوه په~$\sFmla{\True}{!A}$ او بله په~$\sFmla{\False}{!A}$ پيلېږي۔
  
  په کيڼ لوري کښې د لومړي !!{tableau} د فرضونو لاندې برخه ورزياته
  کړئ۔ دلته د قاعدې هره کارونه لا هم سمه ده؛ ځکه د لومړي
  !!{tableau} هر فرض، د~$\sFmla{\True}{!A}$ په شمول، موجود دے۔
  نو د~$\sFmla{\True}{!A}$ لاندې هره څانګه تړل کېږي۔
  
  په ښي لوري کښې د دويم !!{tableau} د فرض لاندې برخه ورزياته کړئ،
  خو د~$\sFmla{\True}{\lnot !A}$ پر فرض د
  $\TRule{\True}{\lnot}$ د هرې کارونې پايلې ترې لرې کړئ۔ پر
  $\sFmla{\True}{\lnot !A}$ د~$\TRule{\True}{\lnot}$ پايله
  $\sFmla{\False}{!A}$ ده، چې بيا هم موجوده ده؛ ځکه د ګډ
  !!{tableau} په ښي لوري کښې د~\Cut{} قاعدې پايله ده۔

  که په دويم تابلو کښې کومه څانګه د فرض
  $\sFmla{\True}{\lnot !A}$ (چې په ګډ !!{tableau} کښې نور فرض نۀ
  دے) او~$\sFmla{\False}{\lnot !A}$ دواړو له امله تړلې وه، پر
  $\sFmla{\False}{\lnot !A}$ د~$\TRule{\False}{\lnot}$ په پلي کولو
  $\sFmla{\True}{!A}$ اخستلے شو۔ اوس په ګډ !!{tableau} کښې اړونده
  څانګه هم تړل کېږي؛ ځکه د~\Cut{} قاعدې ښۍ پايله،
  $\sFmla{\False}{!A}$، پکې ده۔ که په دويم !!{tableau} کښې يوه
  څانګه د بل کوم لامل له امله تړلې وه، په ګډ !!{tableau} کښې
  اړونده څانګه هم تړلې ده؛ ځکه په زوړ، دويم !!{tableau} کښې پر
  څانګه له~$\sFmla{\True}{\lnot !A}$ پرته هر بل !!{signed formula}s
  د ګډ !!{tableau} پر اړوندې څانګې هم راځي۔
  \end{proof}

\end{document}

